\documentclass[10pt,a4paper]{amsart}
\usepackage[margin=19mm]{geometry}
\usepackage{amsmath,amssymb,mathtools}
\usepackage[scr=rsfs,cal=euler]{mathalfa}
\usepackage{xfrac}
\usepackage[unicode,hidelinks,bookmarks=false]{hyperref}
\newcommand{\ie}{i.e.\ }
\newcommand{\cf}{cf.\ }
\newcommand{\resp}{respectively\ }
\newcommand{\Subsection}{\subsection}
\providecommand{\nobreakdash}{\nobreak\hbox{-}}
\setlength{\emergencystretch}{2em}
\multlinegap0pt
\usepackage{amssymb}
\usepackage{enumerate}
\usepackage{mathtools}
\usepackage{tikz}
\usepackage[shortlabels]{enumitem}
\newtheorem{proposition}{Proposition}
\usepackage{cleveref}
\crefname{proposition}{Proposition}{Propositions}
\newcommand\mb{\mathbb}
\newcommand\mc{\mathcal}
\newcommand\mf{\mathfrak}
\title{Moduli Space of Sheaves and Categorified Commutator of Functors}
\author{Yu Zhao}
\date{Source: arXiv:2112.12434v3 (CC BY 4.0)}
\begin{document}
\maketitle

\noindent\emph{Source and licence.} Moduli Space of Sheaves and Categorified Commutator of Functors. Version arXiv:2112.12434v3 (CC BY 4.0), distributed by arXiv under the Creative Commons Attribution 4.0 International licence, http://creativecommons.org/licenses/by/4.0/, as recorded in the arXiv licence metadata for this exact version and shown on the abstract page https://arxiv.org/abs/2112.12434v3. Full licence notice: \texttt{licenses/arxiv-2112.12434-CC-BY-4.0.txt}.\par\smallskip

\noindent\textit{Editorial scope.} Counted unit: the article's proposition lcor (source0.tex lines 826--837). The article's complete original proof of it is delivered with it (source0.tex lines 838--879, one \texttt{proof} environment of 42 lines). No mathematics is written, repaired or re-derived: the statement and the proof are the article's own lines, in the article's own notation, and no retained line is substituted or re-flowed.  Retained uncounted as the source-local context the proof needs: lines 197--204; lines 498--506; lines 531--534; lines 732--735; lines 737--740.  Nothing was omitted from any retained span, so the package carries no omission record.  No retained line contains a citation, so the wrapper carries no bibliography and no bibliography entry.  No retained line contains a figure environment, an image-inclusion command or drawing code, so no image is delivered and the package declares no asset dependency.  Editorial formatting only, in addition: an AMS \texttt{amsart} preamble that restates the article's own declarations and macros used by the retained lines, namely the environments \texttt{proposition} on separate counters, together with the standard packages those lines need.  The retained environments print in this document as Proposition 1 in order of appearance, whereas the article prints the same items with the numbering of its own full document.  The complete native source of the article as distributed in the arXiv e-print for this exact version is bundled unchanged as \texttt{sources/120050/source0.tex}.
  \begin{equation}
\label{hart}
     \mathrm{R}pr_{U*}(\mathcal{O}_{\mb{P}_{X}(U)}(m))\cong  \begin{cases}
      S^{m}U & m\geq 0\\
      0 & -r<m<0 \\
      \mathrm{det}(U)^{-1}\otimes \wedge^{-m-r}(U^{\vee})[1-r]& m\leq -r  
    \end{cases}
  \end{equation}

  \begin{equation}
      \label{cor:5.3ext}
     R(p_{k}\times \pi_{k})_{*}(\mc{L}^{m})=
     \begin{cases}
       S^{m}(\mc{U}_{k}) & m\geq 0 \\
       0 & -r<m<0 \\
     \det(\mc{U}_{k})^{-1}\wedge^{-m-r}(\mc{U}_{k}^{\vee})[1-r] & m\leq -r
     \end{cases}
   \end{equation}

\begin{equation}
  \label{psik}
  \psi_{k,k+1}:\mc{W}_{k}\to \mc{V}_{k+1}
\end{equation}

    \begin{equation}
    \label{eq:neq1}
    h_{i+j-r}^{a+}\cong q^{-a}det(\mc{U}_{k})^{-1}\{\wedge^{r-1+a}\mc{U}_{k}\otimes S^{i+1-a}\mc{U}_{k}\cdots \to S^{i+j}\mc{U}_{k}\}[2-2a-r],
  \end{equation}

  \begin{equation}
    \label{eq:neq2}
    h_{i+j-r}^{a-}\cong q^{-a}det(\mc{U}_{k})^{-1}\{\wedge^{-i-j}\mc{U}_{k}^{\vee}\to \cdots \to \wedge^{-i-j+a}\mc{U}_{k}^{\vee}\otimes S^{-a}\mc{U}_{k}^{\vee}\}[1-a-r]. 
  \end{equation}

\begin{proposition}
  \label{lcor} \label{n54} \label{n56} Equations \cref{eq:neq1} and \cref{eq:neq2} hold. Moreover, $h_{i}^{a+}\cong0$ if $a\leq -r$, and $h_{i}^{a-}\cong0$ if $i<a\leq i+r$ and $i+r\neq 0$. When $i=-r$ and $-r<a\leq 0$, we have
  \begin{equation}
    \label{eq:neq3}
   h_{i}^{a-}\cong q^{-a}det(\mc{U}_{k})^{-1}\mc{O}_{\mc{M}_{k}\times S}[1-2a-r].
 \end{equation}
 Moreover, \cref{eq:neq3} is the image of the truncation morphism $\mf{h}_{-r}^{a-}\to \wedge^{-a}(q\mc{L}\mc{V}_{k+1})^{\vee}[-a]$ under the functor $R((p_{k}\times \pi_{k})\circ pr_{q\mc{L}\otimes \mc{W}_{k}^{\vee}})_{*}$:
 \begin{equation}
   \label{eq:4115}
    h_{i}^{a-}\cong R((p_{k}\times \pi_{k})\circ pr_{q\mc{L}\otimes \mc{W}_{k}^{\vee}})_{*}(\mc{L}^{-r}\wedge^{-a}(q\mc{L}\mc{V}_{k+1})^{\vee}[-a]).
 \end{equation}
\end{proposition}

\begin{proof}
The short exact sequence
  $$0\to \mc{L}^{-1}\mc{V}_{k+1}\to \mc{L}^{-1}\mc{V}_{k}\to \mc{O}_{D}\to 0$$
  on $D$ induces the following long exact sequences on $D$ for any $m>0$:
\begin{align}
  \label{lequa4.5}  0 \to \mc{O} \to \mc{L}\mc{V}_{k}^{\vee} \to \cdots \to \wedge^{m}(\mc{L}\mc{V}_{k}^{\vee})\to \wedge^{m}(\mc{L}\mc{V}_{k+1}^{\vee})\to 0,\\
 \label{lequa4.6} 0\to \wedge^{m}(\mc{L}^{-1}\mc{V}_{k+1})\to \wedge^{m}(\mc{L}^{-1}\mc{V}_{k})\to \cdots \to \mc{L}^{-1}\mc{V}_{k}\to \mc{O}\to 0. 
\end{align}
Recalling the morphism $\mc{L}^{-1}\psi_{k,k+1}:\mc{W}_{k}\to \mc{V}_{k+1}$ from \cref{psik}, \cref{lequa4.5} and \cref{lequa4.6} give
\begin{align}
    \label{eq311}
  S^{m}(\mc{L}\otimes \psi_{k,k+1}^{\vee})&\cong \{\mc{O}\to \cdots \to S^{m-1}(\mc{L}\mc{U}_{k}^{\vee})\to S^{m}(\mc{L}\mc{U}_{k}^{\vee})\}[m],\\
  \label{eq312}
  \wedge^{m}(\mc{L}^{-1}\otimes \psi_{k,k+1})&\cong \{\wedge^{m}(\mc{L}^{-1}\mc{U}_{k})\to \wedge^{m-1}(\mc{L}^{-1}\mc{U}_{k})\to \cdots \to \mc{O}\}[-m].
\end{align}

By \cref{hart}, if $a\leq -r$, then $Rpr_{w_{k}*}\mf{H}^{a+}\cong0$, and hence $h_{i}^{a+}\cong0$. When $a> -r$, \cref{eq312} gives
\begin{align*}
  Rpr_{w_{k}*}\mf{H}^{a+}
                         &\cong det(q\mc{L}\mc{V}_{k+1}^{\vee})det(q\mc{L}\mc{W}_{k}^{\vee})^{-1}\wedge^{rank(\mc{V}_{k+1})-rank(\mc{W}_{k})+a}(q^{-1}\mc{L}^{-1}\psi_{k,k+1})[1-a] \\
                         &\cong q^{r-1}det(\mc{L}\mc{V}_{k+1}^{\vee})det(\mc{L}\mc{W}_{k}^{\vee})^{-1}q^{-r+1-a}\wedge^{r-1+a}(\mc{L}^{-1}\psi_{k,k+1})[1-a] \\
                         &\cong q^{-a}\mc{L}^{r}det(\mc{U}_{k})^{-1}\wedge^{r-1+a}(\mc{L}^{-1}\psi_{k,k+1})[1-a] \\
                       &\cong q^{-a}\mc{L}^{r}det(\mc{U}_{k})^{-1}\{\wedge^{r-1+a}(\mc{L}^{-1}\mc{U}_{k})\to \cdots \to \mc{L}^{-1}\mc{U}_{k}\to \mc{O}\}[2-2a-r]. 
\end{align*}
Thus \cref{eq:neq1} follows from \cref{cor:5.3ext}. When $a\leq 0$, by \cref{hart} and \cref{eq311}, we have
\begin{equation}
  \label{714}
  Rpr_{w_{k}*}\mf{H}^{a-}\cong q^{-a}S^{-a}(\mc{L}\psi_{k,k+1}^{\vee})\cong q^{-a}\{\mc{O}\to \cdots \to S^{-a}(\mc{L}\mc{U}_{k}^{\vee})\}[-a].
\end{equation}
This gives \cref{eq:neq2} and \cref{eq:neq3}. When $i<a\leq i+r$ and $i+r\neq 0$, \cref{714} gives
$$h_{i}^{a-}\cong q^{-a}det(\mc{U}_{k})^{-1}\{\wedge^{i+r}(\mc{U}_{k}^{\vee})\to \cdots \to S^{i+r}(\mc{U}_{k}^{\vee})\}\cong0.$$

It remains to prove \cref{eq:4115}. By \cref{lequa4.5},
  \begin{align*}
    Rpr_{q\mc{L}\mc{W}_{k}^{\vee}}(\mc{L}^{-r}\wedge^{-a}(q\mc{L}\mc{V}_{k+1})^{\vee}[-a])&\cong \mc{L}^{-r}\wedge^{-a}(q\mc{L}\mc{V}_{k+1})^{\vee}[-a] \\
    &\cong q^{-a}\{\mc{O}\to \mc{L}\mc{V}_{k}^{\vee}\to \cdots \to \wedge^{-a}(\mc{L}\mc{V}_{k}^{\vee})\}[-a].
  \end{align*}
The truncation morphism $S^{m}(\mc{L}\mc{U}_{k}^{\vee})\to \wedge^{m}(\mc{L}\mc{V}_{k}^{\vee})$
  induces the morphism
  $$R(pr_{q\mc{L}\mc{W}_{k}^{\vee}})_{*}\mf{h}_{-r}^{a-}\to R(pr_{q\mc{L}\mc{W}_{k}^{\vee}})_{*}(\mc{L}^{-r}\wedge^{-a}(q\mc{L}\mc{V}_{k+1})^{\vee}[-a]).$$
Hence, by \cref{cor:5.3ext} and \cref{n56}, we obtain the isomorphism \cref{eq:4115}.
\end{proof}

\end{document}
